% Person 1 report section for Milestone 2.

\section{Solr Query Configuration: Setup A}

Setup A is the baseline: plain field queries with the default Lucene parser, no
field weighting, no phrase proximity, and no minimum-match tuning. Review-based
needs (Q1, Q2) query the nested \texttt{review} children through Solr's block-join
parent parser with \texttt{score=max}, exactly as in Setup B, so that query
sophistication is the only variable between the two setups rather than the
underlying document model. Plot-based needs (Q3--Q5) search the \texttt{plot} field
directly. The output contract matches Setup B: one row per movie, numeric
\texttt{movie\_id} as the run and qrels identifier, Solr's \texttt{id=movie-<movie\_id>}
retained for traceability.

For Q1 and Q2, the block-join query searches \texttt{review\_text} for a small,
hand-picked set of terms (e.g.\ \textit{clever}, \textit{witty},
\textit{intellectually sharp} for Q1) with plain \texttt{OR} matching; a hard genre
filter restricts Q1 to comedies and Q2 to romance movies. For Q3, an early version
matched \textit{astronaut} and \textit{stranded}/\textit{survive} as independent
\texttt{OR} terms, which over-matched: the plot field alone returned more than
1{,}000 candidates on either word in isolation, so the adopted query requires
\textit{astronaut} \texttt{AND} at least one of \textit{stranded}/\textit{survive}
together in the same document. Q4 similarly requires \textit{FBI} \texttt{AND}
\textit{"serial killer"} jointly, since either concept alone matches many unrelated
crime films. Q5 combines a hard cast-name filter for Robert De Niro with a plain
\texttt{OR} search over mafia vocabulary in the plot field, with an explicit
exclusion for \textit{"Russian mafia"}: an earlier version without this exclusion
returned \textit{Ronin}, confirmed by manual review to be about Russian, not
Italian-American, organized crime.

Setup A returns 162, 91, 14, 7, and 8 documents for Q1--Q5, respectively ---
smaller candidate pools than Setup B's broader, weighted matching, consistent with
Setup A's role as the precision-oriented baseline. Full per-query results and
parameters are saved in \texttt{runs/A\_Q1.json} through \texttt{runs/A\_Q5.json},
and the configuration table and A/B comparison are presented in the following
section; Table~\ref{tab:setup-ab} summarizes the differences between the two
configurations.
